\chapter{外测度}

\begin{solution}{outer-measure-covering-proof}
给定 $A\subseteq\bigcup_{k\ge1}A_k$。若 $\sum_{k\ge1}\mu^*(A_k)=+\infty$，
所需不等式自动成立。否则固定 $\varepsilon>0$，为每个 $k\ge1$ 取覆盖
$A_k\subseteq\bigcup_jC_{k,j}$，使
$\sum_j\ell(C_{k,j})<\mu^*(A_k)+\varepsilon2^{-k}$。双重指标族至多可数，
可排成单列并覆盖 $A$，故
$\mu^*(A)\le\sum_k\mu^*(A_k)+\varepsilon$。令 $\varepsilon\downarrow0$。
\end{solution}

\begin{solution}{outer-null-measurable}
若 $Z\subseteq N$ 且 $\mu^*(N)=0$，则对任意 $A$，
$\mu^*(A\cap Z)=0$。次可加性给出
$\mu^*(A)\le\mu^*(A\cap Z)+\mu^*(A\setminus Z)=\mu^*(A\setminus Z)$，
而单调性给出反向不等式。因此 $Z$ 可测且外测度为零。
\end{solution}

\begin{solution}{caratheodory-finite-unions}
补集封闭由条件的对称性得到。若 $E,F$ 可测，先按 $E$ 分割测试集 $A$，
再按 $F$ 分割 $A\setminus E$，把落在 $E\cup F$ 中的
两项用次可加性合并，得到
$\mu^*(A)\ge\mu^*(A\cap(E\cup F))+\mu^*(A\setminus(E\cup F))$。
另一方向自动成立，故有限并封闭。差写成 $E\cap F^c$。
\end{solution}

\begin{solution}{caratheodory-countable-additivity}
对不交可测集 $E_n$ 与 $E=\bigcup_nE_n$，有限并 $F_N$ 可测，故
\[
 \mu^*(E)=\sum_{n\le N}\mu^*(E_n)+\mu^*(E\setminus F_N)
 \ge\sum_{n\le N}\mu^*(E_n).
\]
令 $N\to\infty$ 得“$\ge$”。外测度的次可加性给出“$\le$”。这也说明
限制函数在可测 $\sigma$-代数上可数可加。
\end{solution}

\begin{solution}{metric-outer-measure-borel}
若 $F=\varnothing$，可测性直接成立。以下设 $F$ 非空且闭，并固定
$A\subseteq X$。若 $\mu^*(A)=+\infty$，
Carathéodory 不等式自动成立。以下设 $\mu^*(A)<\infty$，并令
\[
 A_n=\{x\in A:d(x,F)\ge 1/n\}.
\]
$A_n$ 与 $A\cap F$ 的距离为正，故
\[
 \mu^*(A)\ge\mu^*(A_n\cup(A\cap F))
 =\mu^*(A_n)+\mu^*(A\cap F).
\]
还须证明 $\mu^*(A_n)\to\mu^*(A\setminus F)$；这一步不能诉诸一般外测度
并不具有的从下连续性。令
\[
 C_k=\{x\in A:1/(k+1)\le d(x,F)<1/k\}\qquad(k\ge1).
\]
函数 $x\mapsto d(x,F)$ 为 $1$-Lipschitz；同奇偶性的两个不同指标至少相差2，
对应距离值区间之间留有正间隙。因此任意有限个同奇偶性的非空 $C_k$ 两两
保持正距离，逐次把一块与前面有限块的并合并时也保持正距离。度量可加性与
$\bigcup C_k\subseteq A$ 给出
\[
 \sum_{k\ge1}\mu^*(C_k)\le 2\mu^*(A)<\infty.
\]
另一方面，
\[
 (A\setminus F)\setminus A_n\subseteq\bigcup_{k\ge n}C_k.
\]
由次可加性，右侧的外测度随 $n\to\infty$ 趋于零；所以
$\mu^*(A\setminus F)\le\mu^*(A_n)+o(1)$。结合单调性可得所需极限。
令 $n\to\infty$，便有
\[
 \mu^*(A)\ge\mu^*(A\setminus F)+\mu^*(A\cap F).
\]
故 $F$ 可测。闭集生成 Borel $\sigma$-代数，结论成立。
\end{solution}

\begin{solution}{premeasure-outer-agreement}
自身覆盖给出 $\mu^*(A)\le\mu_0(A)$。对任意代数覆盖 $(A_n)$，令
$B_n=A\cap(A_n\setminus\bigcup_{k<n}A_k)$。预测度在并仍为 $A$ 时可数
可加，且 $B_n\subseteq A_n$，所以
$\mu_0(A)=\sum_n\mu_0(B_n)\le\sum_n\mu_0(A_n)$。取下确界得等号。
再把任意覆盖集按 $A$ 与 $A^c$ 分割，逐项使用有限可加性并取下确界，便得
Carathéodory 条件。
\end{solution}

\begin{solution}{finite-partition-premeasure}
代数中的集合都是若干原子之并。定义
$\mu_0(E)=\sum_{A_j\subseteq E}w_j$。在这个由有限个原子组成的代数上，
不交集合族至多有 $m$ 个非空成员，每个原子只贡献一次，故预测度条件成立；
这里不要求 $X$ 是有限集。

对任意 $S\subseteq X$，每个覆盖 $S$ 的代数集合族必须包含所有与 $S$ 相交
的原子；而这些原子的并本身就是一个允许的覆盖。因此
\[
 \mu^*(S)=\sum_{j:S\cap A_j\ne\varnothing}w_j.
\]
若 $w_j>0$ 且 $E$ 把 $A_j$ 分成两个非空部分，以 $A_j$ 为测试集，
Carathéodory 等式就要求 $w_j=2w_j$，与有限正权重矛盾。因此每个正权重
原子必须全部包含在 $E$ 内或全部在其外。零权重原子的任意子集不受限制。
反之，满足此条件的 $E$ 对任意测试集 $S$，把每个正权重原子的贡献恰好
分到内外两侧之一，零权重块无论怎样分割均贡献零，所以满足可测性等式。

记 $Z=\bigcup_{w_j=0}A_j$。所求可测域及测度为
\[
 \mathcal M_{\mu^*}
 =\left\{\left(\bigcup_{j\in J}A_j\right)\cup N:
 J\subseteq\{j:w_j>0\},\ N\subseteq Z\right\},
\]
对于其中的集合，$\mu((\bigcup_{j\in J}A_j)\cup N)=\sum_{j\in J}w_j$。
这是原有限代数上测度的完备化；若全部权重为零，则可测域是 $\mathcal P(X)$。
\end{solution}

\begin{solution}{extension-nonunique-infinite-total}
记正文的区间代数为 $\mathcal A_0$。每个生成区间都是Borel集，故
$\sigma(\mathcal A_0)\subseteq\mathcal B(\R)$。反过来，每个开集是其中所含的
有理端点半开区间的并：给定开集内一点，可在它两侧取仍留在开集内的有理
端点。这些区间至多可数，因此每个开集都在 $\sigma(\mathcal A_0)$ 中。
两者相等。

计数测度对有限集取其基数，对无限集取无穷，满足可数可加性。
定义 $\nu(E)=0$ 当 $E$ 至多可数，$\nu(E)=+\infty$ 当 $E$ 不可数。
若不交Borel集列每项至多可数，则其并至多可数，等式两端都为零；
若有不可数的一项，则并也不可数，等式两端都为无穷。因此 $\nu$ 是测度。
代数中非空集合都含区间，故两种测度在其上都取无穷，都是 $\mu_0$ 的延拓。
它们在 $\{0\}$ 上分别取 $1$ 和 $0$，确实不同。

相反，在 $\N$ 上设有限集取零、无限集取无穷，则每个单点的值都是零，
而 $\N=\bigcup_{n\ge0}\{n\}$ 的值为无穷。可数可加性要求它等于
$\sum_{n\ge0}0=0$，矛盾。可数个零相加不能因非零项的“个数无穷”而变成无穷。
\end{solution}

\begin{solution}{sigma-finite-disjoint-cover}
令 $Y_1=X_1$，$Y_n=X_n\setminus\bigcup_{k<n}X_k$。代数对有限并和差
封闭，所以 $Y_n\in\mathcal A_0$；它们两两不交、并为 $X$，且
$\mu_0(Y_n)\le\mu_0(X_n)<\infty$。
\end{solution}

\begin{solution}{pi-lambda-theorem}
任意多个 Dynkin 系统的交仍是 Dynkin 系统，且 $\mathcal P(X)$ 包含
$\mathcal P$，故存在最小者 $\lambda(\mathcal P)$，简记为 $\mathcal L$。
先注意 Dynkin 系统内 $U\subseteq V$ 时，
$V\setminus U=(U\cup V^c)^c$ 仍在系统内，因为右边所用的并是不交并。对
$A\in\mathcal P$ 定义
$\mathcal D_A=\{B\subseteq X:A\cap B\in\mathcal L\}$。
它包含 $X$，因为 $A\in\mathcal L$；对补封闭，因为
$A\cap B^c=A\setminus(A\cap B)$ 是 $\mathcal L$ 中的嵌套差集；
对不交可数并封闭，因为 $A\cap\bigcup_nB_n=\bigcup_n(A\cap B_n)$，
且右边仍为不交并。因此 $\mathcal D_A$ 是 Dynkin 系统。
$\mathcal P$ 对交封闭，故 $\mathcal P\subseteq\mathcal D_A$，于是
$\lambda(\mathcal P)\subseteq\mathcal D_A$。再固定
$B\in\lambda(\mathcal P)$，令
$\mathcal E_B=\{A\subseteq X:A\cap B\in\mathcal L\}$。
前一步给出 $\mathcal P\subseteq\mathcal E_B$；它包含 $X$，因为 $B\in\mathcal L$。
又 $A^c\cap B=B\setminus(A\cap B)$ 为嵌套差集，交与不交并的分配律也仍
适用，故 $\mathcal E_B$ 是 Dynkin 系统，包含 $\lambda(\mathcal P)$。所以
$\lambda(\mathcal P)$ 对有限交封闭。一个既是 Dynkin 系统又对有限交封闭
的集合族对补和有限并封闭；再将任意可数并化为不交并，就得到
$\sigma$-代数，故
$\sigma(\mathcal P)=\lambda(\mathcal P)\subseteq\mathcal D$。
\end{solution}

\begin{solution}{outer-measure-from-existing-measure}
空集条件和单调性由定义得到。证明可数次可加性时，只需考虑
\[
 \sum_{n=1}^{\infty}\mu^*(A_n)<\infty
\]
的情形。给定 $\varepsilon>0$，为每个 $n$ 取
$E_n\in\mathcal A$，使
\[
 A_n\subseteq E_n,
 \qquad
 \mu(E_n)<\mu^*(A_n)+\varepsilon2^{-n}.
\]
于是 $\bigcup_nE_n$ 是 $\bigcup_nA_n$ 的可测外包，故
\[
 \mu^*\left(\bigcup_nA_n\right)
 \le \mu\left(\bigcup_nE_n\right)
 \le \sum_n\mu(E_n)
 <\sum_n\mu^*(A_n)+\varepsilon.
\]
令 $\varepsilon\downarrow0$，便知 $\mu^*$ 是外测度。

先证明 $\mathcal A\subseteq\mathcal M_{\mu^*}$。若 $E\in\mathcal A$，
$B\subseteq X$，则对任意可测外包 $H\supseteq B$，
\[
 \mu(H)=\mu(H\cap E)+\mu(H\setminus E)
 \ge \mu^*(B\cap E)+\mu^*(B\setminus E).
\]
对 $H$ 取下确界，再与外测度的次可加性合并，得到 $E$ 的
Carathéodory 等式。外测度零集的任意子集也 Carathéodory 可测，所以原
测度空间的完备化包含在 $\mathcal M_{\mu^*}$ 中。

反过来，先设 $A\in\mathcal M_{\mu^*}$ 且 $\mu^*(A)<\infty$。取可测集
$E_n\supseteq A$，使
$\mu(E_n)<\mu^*(A)+2^{-n}$，并令
\[
 F_n=\bigcap_{k=1}^nE_k,
 \qquad F=\bigcap_{n=1}^{\infty}F_n.
\]
则 $F_n\downarrow F$、$A\subseteq F$，且 $\mu(F_1)<\infty$。由测度的
下降连续性，
\[
 \mu(F)=\lim_n\mu(F_n)=\mu^*(A).
\]
因为 $A$ 满足 Carathéodory 条件，
\[
 \mu(F)=\mu^*(F)
 =\mu^*(F\cap A)+\mu^*(F\setminus A)
 =\mu^*(A)+\mu^*(F\setminus A),
\]
所以 $F\setminus A$ 外测度为零。为回到原空间的完备化，还须找到原可测零包。
一般地，若 $\mu^*(Z)=0$，由下确界定义可取 $H_k\in\mathcal A$，
$Z\subseteq H_k$ 且 $\mu(H_k)<2^{-k}$。于是
$H=\bigcap_{k\ge1}H_k\in\mathcal A$ 包含 $Z$，并且
$0\le\mu(H)\le2^{-k}$ 对所有 $k$ 成立，故 $\mu(H)=0$。
应用于 $Z=F\setminus A$，便知 $A$ 与原可测集 $F$ 的差包含于原可测零集。

最后使用 $\sigma$-有限性。取两两不交的可测集 $X_j$，使
$X=\bigcup_jX_j$ 且 $\mu(X_j)<\infty$。每个 $A\cap X_j$ 仍是
Carathéodory 可测集，且外测度有限；由上一步，存在
$F_j\in\mathcal A$，可令 $F_j\subseteq X_j$，并使
$F_j\setminus(A\cap X_j)$ 包含于某个原可测零集 $H_j$。令 $F=\bigcup_jF_j$，则
$F\in\mathcal A$，而 $F\mathbin\triangle A\subseteq\bigcup_jH_j$，右边是原可测零集。
因此每个 Carathéodory 可测集都属于原测度空间的完备化。这里把全空间分成
有限测度块正是 $\sigma$-有限性的用途。
\end{solution}

\begin{solution}{outer-measure-restriction}
三个外测度公理在子集 $Y$ 上直接继承。若 $E$ 在 $X$ 中
Carathéodory 可测，则对任意 $A\subseteq Y$ 使用同一分割公式，说明
$E\cap Y$ 在 $Y$ 中可测。因此总有
\[
 \{E\cap Y:E\in\mathcal M_{\mu^*}\}\subseteq\mathcal M_{\nu^*}.
\]
若 $Y\in\mathcal M_{\mu^*}$，反向也成立。取 $B\in\mathcal M_{\nu^*}$，
对任意 $A\subseteq X$ 先按 $Y$ 分割，再在 $Y$ 内按 $B$ 分割，得到
\[
 \mu^*(A)=\mu^*(A\cap B)+\mu^*((A\cap Y)\setminus B)+\mu^*(A\setminus Y)
 \ge\mu^*(A\cap B)+\mu^*(A\setminus B).
\]
次可加性给另一方向，故 $B\in\mathcal M_{\mu^*}$，又 $B=B\cap Y$，所求两族相等。

一般 $Y$ 可严格扩大相对可测域。取 $X=\{a,b,c\}$，用
$\varnothing,\{a,b\},\{b,c\}$ 覆盖，空集代价0，另两块各代价1。
诱导外测度在每个单点及 $\{a,b\},\{b,c\}$ 上为1，在 $\{a,c\}$ 及 $X$ 上为2。
三个单点均不可测：分别以 $\{a,b\},\{a,b\},\{b,c\}$ 为测试集，
分割后代价从1变成2。由补对称性，三个二点集也不可测，故
$\mathcal M_{\mu^*}=\{\varnothing,X\}$。但在 $Y=\{a,c\}$ 上，限制外测度
就是计数测度，所有子集均可测，而原可测集的截面只有 $\varnothing,Y$。
\end{solution}

\begin{solution}{premeasure-finite-cover-subadditivity}
令 $B_1=A\cap A_1$，
$B_j=A\cap(A_j\setminus\bigcup_{k<j}A_k)$。则 $B_j$ 两两不交、并为 $A$，
且 $B_j\subseteq A_j$。有限可加性与单调性给出
$\mu_0(A)=\sum_j\mu_0(B_j)\le\sum_j\mu_0(A_j)$。
\end{solution}

\begin{solution}{caratheodory-maximal-domain}
若 $E\in\mathcal C$，任意可测外包 $F\supseteq A$ 都按 $E$ 分解为两个
可测部分，故
$\nu(F)=\nu(F\cap E)+\nu(F\setminus E)$。对外包取下确界，并分别再取
下确界，得到 $E$ 的 Carathéodory 条件。因此
$\mathcal C\subseteq\mathcal M_{\nu^*}$。任意可测零集的子集外测度为零，
自动在 Carathéodory 域中，所以原空间的完备化也包含其中。
\end{solution}
